\section{Conclusion}


Protocol bugs are the hardest bugs to find and the most expensive to exploit. Tamarin
Prover provides machine-checked guarantees that a protocol is secure against all
adversary strategies in the Dolev-Yao model. For Lux Network, Tamarin verification of
the FROST-BLS bridge protocol found a reflection attack that would have caused bridge
liveness failures, and verified that the hybrid PQ handshake provides forward secrecy
even against quantum-capable adversaries.

The investment is modest: each protocol model took 2--3 days to write and seconds to
verify. The return is certainty that the protocol logic is correct, independent of the
implementation.

\begin{thebibliography}{99}
\bibitem{tamarin} Meier, S. et al. ``The TAMARIN Prover for the Symbolic Analysis of Security Protocols.'' CAV, 2013.
\bibitem{dolev1983} Dolev, D. and Yao, A. ``On the Security of Public Key Protocols.'' IEEE Transactions on Information Theory, 1983.
\bibitem{frost} Komlo, C. and Goldberg, I. ``FROST: Flexible Round-Optimized Schnorr Threshold Signatures.'' SAC, 2020.
\bibitem{mlkem} NIST. ``Module-Lattice-Based Key-Encapsulation Mechanism Standard (ML-KEM).'' FIPS 203, 2024.
\bibitem{proverif} Blanchet, B. ``An Efficient Cryptographic Protocol Verifier Based on Prolog Rules.'' CSFW, 2001.
\bibitem{hybrid} Stebila, D. et al. ``Hybrid Key Exchange in TLS 1.3.'' IETF Draft, 2024.
\end{thebibliography}

